% Generated draft: requires mathematical and source review.
\documentclass[11pt]{article}
\usepackage[margin=25mm]{geometry}
\usepackage{fontspec}
\setmainfont{Latin Modern Roman}
\usepackage{amsmath,amssymb,mathtools,unicode-math}
\setmathfont{Latin Modern Math}
\usepackage{enumitem,needspace,xurl,hyperref}
\setlength{\parindent}{0pt}
\setlength{\parskip}{5pt}
\setlength{\emergencystretch}{3em}
\newcommand{\N}{\mathbb{N}}
\newcommand{\Z}{\mathbb{Z}}
\newcommand{\Q}{\mathbb{Q}}
\newcommand{\R}{\mathbb{R}}
\newcommand{\C}{\mathbb{C}}
\newcommand{\F}{\mathbb{F}}
\newcommand{\A}{\mathbb{A}}
\newcommand{\PP}{\mathbb{P}}
\newcommand{\E}{\mathbb{E}}
\newcommand{\eps}{\varepsilon}
\newcommand{\dd}{\,\mathrm{d}}
\begin{document}
\section*{UnsolvedMath ID: 1900008}
\textbf{Source problem code:} catalogue entry
\subsection{Definitions and mathematical statement}
A combinatorial type of a convex polytope is its face lattice, up to an inclusion-preserving isomorphism; this fixes incidence relations rather than a particular geometric realization. In a conventional lattice-sail construction, for a pointed convex cone $C\subset\mathbb R^d$ with apex $0$ and a lattice $\Lambda$, one first forms
\[
K(C,\Lambda)=\operatorname{conv}\big((C\cap\Lambda)\setminus\{0\}\big)
\]
and obtains the sail from the boundary of this lattice convex hull, with the cone-boundary and irrationality conventions fixed by the chosen sail definition. This describes the standard framework only: the supplied record does not specify those conventions or which cones are assigned to a polytope.

\textbf{Multidimensional IKEA problem.} For every prescribed polytope combinatorial type $\tau$, describe all collections of sails of cones that arise from polytopes $P$ whose face lattice is isomorphic to $\tau$. In schematic notation, if $\mathcal S(P)$ denotes the source's intended collection, determine
\[
\{\mathcal S(P):P\text{ is a polytope of combinatorial type }\tau\}.
\]
A precise classification requires recovering the cone assignment, lattice, admissible realizations, dimension scope, and equivalence relation on collections; they are absent from the supplied statement. In particular, tangent cones and normal cones are not silently identified. The supplied research summary reports restricted constructive results and no verified full resolution, with moderate-confidence OPEN-TRIAGE status.
\subsection{Short English statement}
For a fixed polytope incidence pattern, classify the collections of lattice sails realizable by its geometric realizations. The source omits the cone and sail conventions needed to turn this request into a complete classification statement.
\subsection{Sources}
\begin{itemize}
\item O. Karpenkov, \emph{Open problems in geometry of continued fractions} (2017), Problem 8, \url{https://arxiv.org/abs/1712.01450}; source-TeX extraction, accessed 2026-07-29.
\item The supplied research summary refers generically to Karpenkov's papers and monograph for restricted three-dimensional and special-family results, with no identifiable title or theorem citation. It cross-refers to a two-dimensional analogue using the literal code AMS-018-0002; no permanent ID is supplied for that reference.
\item Local registry \texttt{data/unsolvedmath/problems.json}, permanent ID $1900008$, source code catalogue entry, \emph{Geometry of Continued Fractions --- Combinatorial structure of sails}. The record labels this open and OPEN-TRIAGE, with NEEDS\_REVIEW source-status and rights qualifications.
\end{itemize}
\end{document}
